% Source: source/普通高考/2012/2012江西理.pdf, page 1, question 14.
% pdfimages -list found no embedded bitmap; this program flowchart is vector artwork.
% Grid evidence: tmp/review_2012_jiangxi_science/grid/page1_grid100.jpg and
% q14_grid50.jpg; no-grid source crop:
% tmp/review_2012_jiangxi_science/crops/q14_original_full.png.
% Node -> directed-edge list: 开始 -> T=0,k=1 ->
% sin(k*pi/2)>sin((k-1)*pi/2)?; 是 -> a=1, 否 -> a=0, and both branches
% merge before T=T+a; then T=T+a -> k=k+1 -> k<6?; 是 -> left loop-back
% to the shared stem before the first test, 否 -> 输出 T -> 结束.
% All borders and connectors are solid; node text is locally zero-indented and
% centered.  The two branch merges are explicit and no other connectors occur.

\begin{tikzpicture}[
    >=Stealth,
    x=1cm,
    y=0.92cm,
    line width=0.45pt,
    line cap=round,
    line join=round,
    font=\normalsize,
    flowtext/.style={anchor=center,align=center,
        execute at begin node={\setlength{\parindent}{0pt}}},
    every node/.style={flowtext},
    flowbox/.style={draw,flowtext,inner xsep=4pt,inner ysep=2pt,
        minimum height=0.72cm},
    terminal/.style={flowbox,rounded corners=0.18cm,minimum width=1.55cm},
    process/.style={flowbox},
    decision/.style={flowbox,diamond,inner sep=2pt},
    branchlabel/.style={flowtext,inner sep=0pt},
]
    \path[use as bounding box] (-5.00,-1.85) rectangle (6.35,15.40);

    \node[terminal] (start) at (0,14.95) {开始};
    \node[process,minimum width=3.55cm] (init) at (0,13.25) {$T=0,k=1$};
    \node[decision,minimum width=6.20cm,minimum height=1.95cm,aspect=3.0,
        text width=5.20cm]
        (first) at (0,10.75)
        {$\sin\dfrac{k\pi}{2}>\sin\dfrac{(k-1)\pi}{2}?$};

    \node[process,minimum width=2.15cm] (aone) at (0,8.35) {$a=1$};
    \node[process,minimum width=2.15cm] (azero) at (5.0,8.35) {$a=0$};
    \node[process,minimum width=3.75cm] (add) at (0,6.55) {$T=T+a$};
    \node[process,minimum width=3.05cm] (inc) at (0,4.65) {$k=k+1$};
    \node[decision,minimum width=3.35cm,minimum height=1.15cm,aspect=3.2]
        (second) at (0,2.45) {$k<6?$};
    \node[flowbox,trapezium,trapezium left angle=75,trapezium right angle=105,
        minimum width=2.55cm,minimum height=0.86cm]
        (output) at (0,0.35) {输出 $T$};
    \node[terminal] (finish) at (0,-1.25) {结束};

    % Initial stem and first test.
    \draw[->] (start.south) -- (init.north);
    \draw[->] (init.south) -- (first.north);

    % First test: affirmative branch descends to a=1; negative branch goes
    % right to a=0.  Both branches merge immediately above T=T+a.
    \draw[->] (first.south) -- node[branchlabel,right=3pt] {是} (aone.north);
    \draw (aone.south) -- (0,7.30);
    \draw[->] (first.east) -- node[branchlabel,above=3pt] {否} (5.0,10.75)
        -- (azero.north);
    \coordinate (branchMerge) at (0,7.30);
    \coordinate (branchTip) at (0.35,7.30);
    \draw (azero.south) -- (5.0,7.30);
    \draw[->] (5.0,7.30) -- (branchTip);
    \draw (branchTip) -- (branchMerge);
    \draw[->] (branchMerge) -- (add.north);

    % Common update path and second test.
    \draw[->] (add.south) -- (inc.north);
    \draw[->] (inc.south) -- (second.north);
    \draw[->] (second.south) -- node[branchlabel,right=3pt] {否} (output.north);
    \draw[->] (output.south) -- (finish.north);

    % The affirmative second-test branch loops around the left side and
    % re-enters the stem before the first test with one right-pointing arrow.
    \coordinate (loopLeft) at (-4.45,2.45);
    \coordinate (loopTop) at (-4.45,11.85);
    \coordinate (loopArrow) at (-1.20,11.85);
    \draw (second.west) -- node[branchlabel,above=3pt] {是} (loopLeft) -- (loopTop);
    \draw[->] (loopTop) -- (loopArrow);
    \draw (loopArrow) -- (0,11.85);
\end{tikzpicture}
